\documentclass[english, a4paper]{article}

%for writing
\usepackage{lmodern}
\usepackage[utf8]{inputenc}
\usepackage[T1]{fontenc}
\usepackage{babel}

%for layout
\usepackage[
    a4paper,
    left=25mm,
    right=25mm,
    top=25mm,
    bottom=25mm
    ]{geometry}

% for maths
\usepackage{amsmath}
\usepackage{amsfonts}
\usepackage{amssymb}
\usepackage{amsthm}

%for graphics
\usepackage{graphicx}

% for refs (in that order)
\usepackage{hyperref}

% for headers and footers
\usepackage{fancyhdr}

% Custom commands
\usepackage{import}

\newcommand{\keywords}[1]
{
    \small
    \textbf{Keywords:} #1
    \vspace{0.2cm}
}

%document properties
\title{Draft of trapeze like congestion}

\author{Martin Briand$^{a,}$\footnote{Martin Briand is the corresponding author. His corresponding email is \href{mailto:martin.briand23@gmail.com}{martin.briand23@gmail.com}}, Nicolas Coulombel$^{b}$, Guillaume Monchambert$^{c}$\\
\small $^{a}$Martin’s affiliation\\
\small $^{b}$Nicolas’ affiliation\\
\small $^{c}$Guillaume’s affiliation}

\date{2025}

% document
\begin{document}

\maketitle

\begin{abstract}
    No abstract
\end{abstract}

\keywords{No keywords for the moment…}

\section{First idea of a problem setting}

\subsection{Essence of the idea}

The idea is to share the capacity of a highway bottleneck between Trapeze like reserved lanes and casual general purpose lanes. Commuters can choose between 3 modes: driving alone on general purpose lanes, driving on Trapeze lanes while assuming the risk of having to take passengers if asked by the Trapeze system, or being a carpooling passenger on the Trapeze lanes.

When making the decision to be a Trapeze driver, commuters are compensated (potentially negatively), and passengers pay a fee (potentially negative).

\subsection{Formal setting}

\end{document}